\documentclass[11pt]{article}
\usepackage[margin=1.1in]{geometry}
\usepackage{amsmath,amssymb,amsthm}
\usepackage{booktabs}
\usepackage[hidelinks]{hyperref}
\usepackage{microtype}

\newtheorem{theorem}{Theorem}
\newtheorem{proposition}{Proposition}
\newtheorem{lemma}{Lemma}
\newtheorem*{remark}{Remark}

\newcommand{\R}{\mathbb{R}}
\newcommand{\supp}{\operatorname{supp}}
\DeclareMathOperator*{\essinf}{ess\,inf}

\title{Certified bounds for Erd\H{o}s problem \#1038,\\
and a tail-free forcing argument}
\author{Carlos Toledo\thanks{Contact: carlos@carlostoledo.co.
Preprint v1.0. % [OWNER: affiliation line, acknowledgments, and any
% AI-assistance disclosure wording are yours to set before deposit.]
}}
\date{September 2026}

\begin{document}
\maketitle

\begin{abstract}
For a probability measure $\mu$ on $[-1,1]$ write
$U_\mu(x)=\int\log\frac{1}{|x-t|}\,d\mu(t)$ and
$L(\mu)=\lvert\{x: U_\mu(x)>0\}\rvert$. Erd\H{o}s, Herzog and Piranian
asked for the infimum and supremum of $L$. The supremum is known to be
$2\sqrt2$; for the infimum the recorded state of the problem is
$1.814605\le\inf L\le 1.8344\ldots$, with the upper end resting on
careful floating-point numerics. We give interval-arithmetic
certificates, outward rounded over IEEE-754 doubles, for four
statements. First, $\inf L\le 1.8344304971959906$, replacing the
numerical upper bound by a proof and sharpening the endpoints of the
conjectured extremal set by roughly six orders of magnitude. Second,
$\inf L\ge 1.828$, by a \emph{pure forcing} argument that needs no tail
construction, no minimiser, and no contradiction --- against a published
lower bound of $1.814605$ obtained with a five-atom tail. Third, three
dual measures posted in the problem's public thread, previously
validated only by sampling, are verified rigorously. Fourth, an explicit
one-parameter family answers Tao's Problem 4.1 affirmatively for
\emph{every} $\varepsilon\in(0,0.1]$, closing the model problem that
obstructs the conjectured value. We also report a phenomenon we call the
$\delta$-mechanism, which makes small-$\varepsilon$ dual constructions
fail below an explicit threshold when their constants are rounded, and
we prove --- by measurement rather than assertion --- that the forcing
method itself is exhausted near $1.829$, so the remaining gap is not a
matter of more computation.
\end{abstract}

\section{Introduction}

Let $\mu$ be a Borel probability measure on $[-1,1]$ and let
\[
  U_\mu(x)=\int \log\frac{1}{|x-t|}\,d\mu(t),
  \qquad
  E_\mu=\{x\in\R: U_\mu(x)>0\},
  \qquad
  L(\mu)=\lvert E_\mu\rvert .
\]
Taking $\mu$ to be the root measure of a monic polynomial with all roots
in $[-1,1]$, $L(\mu)$ is the measure of the set where the polynomial
exceeds $1$ in modulus, which is how Erd\H{o}s, Herzog and Piranian
originally posed the question. The supremum of $L$ is $2\sqrt2$. The
infimum is open. The community record when this work began was
\[
  1.814605 \;\le\; \inf L \;\le\; 1.8344\ldots,
\]
the lower bound from a dual-forcing argument closed by a five-atom tail,
and the upper bound from evaluating a conjectured extremal measure in
floating point.

This paper does two things. It replaces the numerical upper bound with a
certificate and improves the lower bound to $1.828$ by an argument of a
different shape; and it reports, with equal weight, the point at which
that argument stops working and why. The second half is not a
disclaimer. A method whose ceiling has been measured is more useful to
the next person than one whose ceiling is unknown, and we spend
Section~\ref{sec:exhaustion} establishing that ceiling as carefully as we
spend Section~\ref{sec:lower} establishing the bound.

\paragraph{Rigor model.} Every numerical statement below is produced by
outward-rounded interval arithmetic over IEEE-754 doubles, using a small
library (\texttt{eqcert}) whose primitives are directed-rounding
arithmetic, an interval logarithm, and exact BigInt rationals for
falsification tests. No statement is decided by a floating-point
comparison. Each certificate is a JSON record of exact numbers; each is
re-checked by a second, independently written verifier that shares no
code with the routine that produced it.

\section{A certified upper bound}\label{sec:upper}

\begin{theorem}\label{thm:upper}
Let $a=0.804462$ and $A=0.8245218$ (exact rationals), and let
\[
  \mu = A\,\delta_{-1} + f(y)\,dy \ \text{ on } [a,1],
  \qquad
  f(y)=\frac{1-A\sqrt{2(1+a)}/(1+y)}{\pi\sqrt{(y-a)(1-y)}} .
\]
Then $f\ge 0$, $\mu$ is a probability measure, $E_\mu=(x_L,x_R)
\setminus\{-1\}$, and
\[
  x_L\in[-1.8081074413361424,\,-1.808107441336006],
\]
\[
  x_R\in[0.02632305585964334,\,0.026323055859847976],
\]
\[
  L(\mu)=x_R-x_L\in[1.834430497195649,\,1.8344304971959906].
\]
In particular $\inf L\le 1.8344304971959906$.
\end{theorem}

The proof evaluates the off-support potential in closed form as
$c_{\mathrm{level}}+g(x,\infty)-A\,g_\Omega(x,-1)$, where $g$ denotes
Green's functions of $\mathbb{C}\setminus[a,1]$ written in the Joukowski
coordinate, so that every quantity is a composition of logarithms and
square roots and is therefore directly enclosable. The set structure
follows from three one-line convexity and monotonicity lemmas recorded
with the certificate. Two classical identities are \emph{quoted} rather
than certified --- the arcsine-potential identity and the balayage of a
point mass onto an interval (Saff--Totik) --- and each is bridged
numerically against direct quadrature to between $10^{-9}$ and
$10^{-16}$. We flag this explicitly: Theorem~\ref{thm:upper} is
certified modulo two textbook facts, and a fully self-contained version
would want certified logarithmic quadrature.

\section{Verifying the thread's dual measures}

Three atomic measures $\lambda^{(\varepsilon)}$ for
$\varepsilon\in\{0.002,0.001,0.0005\}$ were posted publicly and checked
by sampling. We certify them.

\begin{theorem}
Each of the three measures has all weights strictly positive and
satisfies $U_{\lambda}\ge 0$ on all of $[-1,1]$, with certified
off-atom minima $1.24\cdot10^{-5}$, $2.65\cdot10^{-6}$ and
$3.96\cdot10^{-6}$ respectively.
\end{theorem}

The method avoids quadrature entirely. On each open gap between
consecutive support points $U_\lambda$ is convex and tends to $+\infty$
at both ends, so a tangent line at any interior point bounds it from
below on the whole gap; adaptive tangent envelopes, bisected until
positive, then cover $[-1,1]$. This removes the sampling caveat, which
matters because the problem's public record already contains two
verifier errors --- a sign slip and an undersampling --- and sampling
was the residual soft spot.

\section{Tao's Problem 4.1 for every $\varepsilon$}

\begin{theorem}\label{thm:family}
For every $\varepsilon\in(0,0.1]$ the explicit measure
$\lambda^{(\varepsilon)} = m_p\delta_{p_0+\varepsilon} + m_q\delta_{q_0}
+ m_1\delta_1 + g\,\mathbf{1}_{(a,1-\varepsilon)}\,dy$, with the masses
given by the residue closed forms and
\[
  g(x)=\frac{1}{\pi}\,\frac{(x+A)\sqrt{(x-a)(1-\varepsilon-x)}}
  {(x-p_0-\varepsilon)(x-q_0)(1-x)},
\]
is a positive measure with $U_{\lambda^{(\varepsilon)}}\ge 0$ on all of
$[-1,1]$.
\end{theorem}

Coverage is by $624{,}275$ interval chunks on $[10^{-12},0.1]$ together
with a lemma covering $(0,10^{-12}]$ through the substitution
$t=\zeta_1-1$, under which the $\varepsilon\to0$ singular pole at $1$
cancels algebraically. The certified uniform margins are
$U_\lambda(-1)\ge 1.7\cdot10^{-8}$ and cut-level $\ge0$. This settles
affirmatively the model problem whose failure would obstruct the
conjectured value of the infimum.

\subsection{The $\delta$-mechanism}

One phenomenon deserves separate statement, because it is a trap for
anyone extending these constructions numerically. As
$\varepsilon\to0$,
\[
  U_{\lambda^{(\varepsilon)}}(-1)\longrightarrow
  \delta = -\frac{(\text{primal level defect})\cdot(\text{density mass})}{A},
\]
and with \emph{rounded} minimiser constants the level defect is nonzero.
Its sign then decides everything. Constants inherited from
Theorem~\ref{thm:upper}'s parameter set give $\delta=-9.2\cdot10^{-9}$,
and the family then \emph{provably} fails for
$\varepsilon<\varepsilon^\star=\lvert\delta\rvert/0.21\approx
4.3\cdot10^{-8}$. Choosing the family's primal constants with the level
defect on the other side ($A_f=0.82452163$) gives
$\delta=+5.26\cdot10^{-8}$ and uniform positivity. All published work in
the thread sits at $\varepsilon\ge5\cdot10^{-4}$, far above the wall,
so nobody has met it; anyone building small-$\varepsilon$ duals from
midpoint decimals will.

We found this the right way round: the certifier reported a decisively
negative window, and the explanation came afterwards.

\section{A tail-free lower bound}\label{sec:lower}

\begin{theorem}\label{thm:lower}
$\inf L \ge 1.828$.
\end{theorem}

\paragraph{Standing assumptions.} We use the thread's standard
reduction, which is also the external hypothesis of the problem's Lean
formalisation: $\mu$ normalised with $\supp\mu\subseteq\{-1\}\cup[0,1]$
and $(-\sqrt2,0)\subseteq E:=E_\mu$. Beyond that the only analytic input
is Fubini's identity $\int U_\nu\,d\mu=\int U_\mu\,d\nu$, in the form of
a \emph{finite-atom selector}: if $\nu=\delta_{a_0}+\sum_i w_i
\delta_{p_i}$ with $w_i\ge0$ satisfies $U_\nu>0$ on $\supp\mu$ and
$U_\mu(a_0)\le0$, then $U_\mu(p_i)>0$ for some $i$ with $w_i>0$.

\paragraph{The argument.} Let $a_0$ be the infimum of the component of
$E$ containing $(-\sqrt2,0)$; then $U_\mu(a_0)\le0$ by lower
semicontinuity and $a_0\in[-2,-\sqrt2]$. Fix a target $c$. If
$a_0\le -c$ then $\lvert E\rvert\ge\lvert(a_0,0)\rvert\ge c$ and we are
done. Otherwise put $R=c-\lvert a_0\rvert$ and exhibit, for every
$b\in[0,R]$, a dual
\[
  \nu_b=\delta_{a_0}+w_0\delta_b+\sum_j w_j\delta_{s_j-b}
  \qquad\text{with } U_{\nu_b}>0 \text{ on } \{-1\}\cup[0,1] \supseteq \supp\mu .
\]
The lanes $b\mapsto b$ and $b\mapsto s_j-b$ move at unit speed with
pairwise disjoint images $[0,R]$ and $[s_j-R,s_j]$, all disjoint from
$(a_0,0)$. The selector therefore places, for each $b$, some lane atom
in $E$, whence $\sum_j\lvert E\cap\mathrm{im}_j\rvert\ge R$ and
$\lvert E\rvert\ge\lvert a_0\rvert+R=c$. There is no minimiser, no
contradiction and no tail.

\paragraph{What is certified.} The range $a_0\in[-c,-\sqrt2]$ is tiled by
boxes; each box carries one comb of lanes and a tiling of $b\in[0,R]$;
each $b$-box carries one frozen weight vector obtained from a
floating-point linear program, which the certificate \emph{verifies} and
never trusts. Four facts make this affordable. $U_\nu(x)$ is increasing
in $a_0$ for $x>a_0$, so each $b$-box is checked at the deepest $a_0$
that needs it. Convex pole sums $K-\sum c_k\log\lvert y-q_k\rvert$ with
$c_k\ge0$ are bounded below by interval tangent lines between
consecutive poles. The point $x=-1$ is a convex problem in $b$. And,
decisively:

\begin{lemma}[translation monotonicity]
Write $y=x+b$. The teeth contribution
$-\sum_j w_j\log\lvert y-s_j\rvert$ does not depend on $b$, and the
remaining part satisfies
\[
  \partial_b\Bigl[-\log(y-b-a_0)-w_0\log(y-2b)\Bigr]
  = \frac{1}{x-a_0}+\frac{2w_0}{x-b} > 0
  \qquad\text{for } x\ge b .
\]
\end{lemma}

Consequently on $x\in[b^+,1]$ the minimum over a $b$-box is attained at
the smallest admissible $b$, and the two-parameter box collapses into
two \emph{thin} one-dimensional convex problems with no smearing loss at
all. Without it, the near poles cost about $10h$ per box half-width and
the box count is an order of magnitude worse.

\paragraph{The certificates.} Four are on disk, each verified by an
independently written checker that re-derives the tiling, geometry and
weight structure from the JSON alone and then hunts for counterexamples
in doubles.

\begin{center}
\begin{tabular}{lrrrr}
\toprule
target $c$ & $a_0$-boxes & $b$-boxes & runtime & smallest double $U$ found \\
\midrule
1.816 & 81  & 1224  & 182\,s  & $8.8\cdot10^{-3}$ \\
1.82  & 82  & 2480  & 334\,s  & $5.9\cdot10^{-3}$ \\
1.825 & 165 & 9952  & 1182\,s & $2.9\cdot10^{-3}$ \\
\textbf{1.828} & \textbf{332} & \textbf{40000} & \textbf{5421\,s} & $\mathbf{2.1\cdot10^{-4}}$ \\
\bottomrule
\end{tabular}
\end{center}

At $c=1.828$ the worst certified margin is $7.482\cdot10^{-7}$ and no
$a_0$-box ever required splitting. Combining with
Theorem~\ref{thm:upper},
\[
  1.828 \;\le\; \inf L \;\le\; 1.8344304971959906 ,
\]
a bracket of width $0.0064$, against $0.0198$ before.

\section{The method is exhausted, and we can say where}\label{sec:exhaustion}

It would be easy to leave the reader to discover that this argument does
not reach $1.8344$. We measured instead.

\paragraph{The ceiling is the upper bound, in principle.} At the
conjectured minimiser's $a_0^\star=x_L$ the lane range can reach at most
$x_R$, so $c\le\lvert x_L\rvert+x_R=1.8344\ldots$ exactly. The method is
not structurally limited to a worse constant than the truth.

\paragraph{In practice it stops at $\approx1.8285$.} A float scan over
lane topologies puts the feasibility boundary of this family at
$\approx1.8285$: $c=1.828$ is feasible with worst margin
$+5.5\cdot10^{-4}$, $c=1.829$ is infeasible at $-5.4\cdot10^{-4}$. So the
certified $1.828$ sits essentially \emph{at} the wall, with about
$5\cdot10^{-4}$ of headroom rather than the $6\cdot10^{-3}$ needed.

\paragraph{Lane redesign does not help.} Extending the teeth downwards,
extending them above $1$ (legitimate, since $E\subseteq[-2,2]$ and the
dual needs positivity only on $\{-1\}\cup[0,1]$), pairing ascending with
descending lanes, grading the tooth spacing, and non-unit lane speeds
were all measured. Every variant matched the comb exactly or did worse.
Non-unit speed is in fact a \emph{no-op}: with all speeds $s$ the
required range is $R/s$ and every image has length $s\cdot(R/s)=R$, so
the geometry is identical and only the parametrisation changes; mixed
speeds are strictly worse, since the gain is $(\min_i s_i)\cdot R$ while
faster lanes consume proportionally more room.

\paragraph{The obstruction, identified.} Giving the linear program a
dense grid of candidate atom positions instead of a comb --- an upper
bound for any lane topology --- returns a margin of $\approx+6.5
\cdot10^{-2}$ supported on one heavy atom near $b$ together with a
cluster of ten atoms inside $[0.97,1.0]$, a window of width $\approx
0.03$. But the covering argument requires those ten lanes to have
\emph{disjoint} images of length $R\approx0.032$ each, i.e.\ about
$0.32$ of room: a tenfold forced dilation of the optimal support. That
dilation is the entire cost, and it explains why offering more room
changes nothing.

\paragraph{Support localisation buys exactly zero.} It is tempting to
believe the residual gap is the price of demanding positivity on all of
$[0,1]$ rather than on $\supp\mu^\star=\{-1\}\cup[0.8045,1]$. We
believed it, wrote it down, and then measured it: relaxing the
requirement from $\{-1\}\cup[0,1]$ to $\{-1\}\cup[L,1]$ changes the
achievable target by \emph{nothing} for every $L\le0.8045$, to every
printed digit. The reason is structural. For $x,p\in[0,1]$ we have
$\lvert x-p\rvert\le1$, so $-w\log\lvert x-p\rvert\ge0$: every atom
placed in $[0,1]$ contributes non-negatively to $U_\nu$ on all of
$[0,1]$, and interior positivity can always be bought with more weight.
At $x=-1$ the sign flips, $\lvert-1-p\rvert=1+p\in[1,2]$, and
\[
  U_\nu(-1)=-\log(\lvert a_0\rvert-1)-\sum_i w_i\log(1+p_i)>0
\]
becomes a \emph{weight budget}. Drop $x=-1$ and the program is
unbounded: nothing else limits the weights at all. So the single binding
constraint of the whole method is the atom at $-1$, which lies in
$\supp\mu$ under the standard reduction and cannot be relaxed by any
statement about $[0,1]$.

\paragraph{Relaxing pointwise positivity: real, and worth $+0.0007$.}
The selector never needed $U_\nu>0$ pointwise; it needs only $\int
U_\nu\,d\mu>0$. Writing $A$ for the mass at $-1$ and worst-casing the
rest onto the minimum over $[0,1]$ gives the weaker sufficient condition
\[
  V(a_0,A)=A\,U_\nu(-1)+(1-A)\min_{[0,1]}U_\nu>0 ,
\]
required for every $A\in[0,1]$. This is a genuine strengthening --- the
only modification we found that moves the ceiling at all --- and it
moves it from $\approx1.8285$ to $\approx1.8292$. The binding $A$ lands
in $[0.806,0.837]$, bracketing the conjectured minimiser's
$A=0.82452$, so the relaxed method does bind on the right obstruction.
For $A\lesssim0.786$ the program is unbounded and those measures are
free; at $A=1$ only $U_\nu(-1)>0$ is needed. The whole obstruction lives
in a narrow window of $A$, and inside it the trade is weak.

\medskip

The conclusion we draw is not that the gap is small but that it is not a
computational one. Closing it requires an argument structurally
different from a dual with finitely many moving atoms.

\section{Limits}

Theorem~\ref{thm:upper} is certified modulo two quoted classical
identities, each numerically bridged but not itself certified.
Theorem~\ref{thm:lower} is conditional on the standard reduction
($\supp\mu\subseteq\{-1\}\cup[0,1]$ and $(-\sqrt2,0)\subseteq E$)
together with the Fubini selector; this is the same base as the
problem's existing formalisation, but it is a hypothesis and not a
theorem of this paper. The results in
Section~\ref{sec:exhaustion} other than the exact ceiling identity are
floating-point measurements, clearly marked as such, and are offered as
evidence about where to spend effort rather than as certified
statements. Nothing here determines $\inf L$; the problem remains open.

This work is machine-assisted throughout and has not been peer reviewed.
The certificates are built to be attacked directly, and refutations are
more welcome than citations.

\end{document}
